\section{References}

\subsection*{Primary Philosophical and Theological Sources}

\begin{itemize}
\item \textbf{Aristotle.} \emph{Posterior Analytics}. Translated by G.\ R.\ G.\ Mure. In \emph{The Works of Aristotle Translated into English}, vol.\ I, edited by W.\ D.\ Ross. Oxford: Clarendon Press, 1928. (Public domain; cited by Bekker pagination: I.1--4, I.6--8, I.13--14, I.30, II.8, II.11, II.16, II.19.)
\item \textbf{Aristotle.} \emph{Prior Analytics}. Translated by A.\ J.\ Jenkinson. Oxford: Clarendon Press, 1928. (Public domain; cited at I.1, 24b18--20.)
\item \textbf{Aristotle.} \emph{Topics}. Translated by W.\ A.\ Pickard-Cambridge. Oxford: Clarendon Press, 1928. (Public domain; cited at I.1--2, 100a18--b23, 101a25--37.)
\item \textbf{Aristotle.} \emph{Metaphysics}. Translated by W.\ D.\ Ross. Oxford: Clarendon Press, 1924. (Public domain; cited at III.3, 998b22; IV [Gamma] 1--4, 1003a21--1009a5.)
\item \textbf{Aristotle.} \emph{Nicomachean Ethics}. Translated by W.\ D.\ Ross. Oxford: Clarendon Press, 1925. (Public domain; cited at VI.3, 1139b18--22; VI.6, 1140b31--1141a8.)
\item \textbf{Duns Scotus, Bl.\ John.} \emph{Ordinatio}, Book I, dist.\ 3, part 1, qq.\ 1--2, n.\ 26. Critical Vatican edition (\emph{Opera Omnia}, vol.\ III). Vatican City: Typis Polyglottis Vaticanis, 1954. (Historical Latin text defining a univocal concept.)
\item \textbf{Hume, David.} \emph{Dialogues Concerning Natural Religion}, Part IX. London, 1779. (Public domain; Cleanthes on the twenty particles of matter.)
\item \textbf{John of Damascus, St.} \emph{De Fide Orthodoxa} (\emph{An Exact Exposition of the Orthodox Faith}), Book I, cc.\ 1, 3. Translated by S.\ D.\ F.\ Salmond, 1899. (Public domain.)
\item \textbf{Thomas Aquinas, St.} \emph{Summa Theologiae}, Prima Pars (qq.\ 1--49). Leonine critical edition, vol.\ IV. Rome: Typographia Polyglotta, 1888 (checked via Corpus Thomisticum). English translation: Fathers of the English Dominican Province, Benziger Brothers, New York (public domain, Project Gutenberg eBook 17611).
\item \textbf{Thomas Aquinas, St.} \emph{Summa contra Gentiles}, Book I. Leonine manual edition. Turin--Rome: Marietti, 1961 (checked via Corpus Thomisticum; cited at cc.\ 3, 10--11, 13 nn.\ 14, 15, 33, 14--28).
\item \textbf{Thomas Aquinas, St.} \emph{De ente et essentia}. Edited by L.\ Baur, emended by J.\ Koch. M\"unster: Aschendorff, 1933 (checked via Corpus Thomisticum; cited at cc.\ 1, 4).
\item \textbf{Thomas Aquinas, St.} \emph{Expositio libri Posteriorum Analyticorum}. Leonine critical edition, vol.\ I. Rome: Typographia Polyglotta, 1882 (checked via Corpus Thomisticum; cited at I, lect.\ 2; II, lect.\ 1, 7, 9).
\item \textbf{Thomas Aquinas, St.} \emph{Expositio super librum Boethii De Trinitate}. Edited by B.\ Decker. Leiden: Brill, 1959 (checked via Corpus Thomisticum; cited at q.\ 5, a.\ 3; q.\ 6, a.\ 4 ad 2).
\end{itemize}

\subsection*{Historical Commentators and Modern Secondary Literature Cited for Open Disputes}

\begin{itemize}
\item \textbf{Cajetan, Thomas de Vio.} \emph{De nominum analogia} (1498). Cited historically as the classical source formulating the syllogistic four-term problem for analogical middle terms.
\item \textbf{Feser, Edward.} \emph{Five Proofs of the Existence of God}. San Francisco: Ignatius Press, 2017. (Cited in Lesson 8 as a modern sympathetic defence of classical proofs.)
\item \textbf{Kenny, Anthony.} \emph{The Five Ways: St.\ Thomas Aquinas' Proofs of God's Existence}. London: Routledge \& Kegan Paul, 1969. (Cited in Lesson 8 as a modern critical assessment arguing that the Five Ways fail as strict demonstrations.)
\item \textbf{Wippel, John F.} \emph{The Metaphysical Thought of Thomas Aquinas: From Finite Being to Uncreated Being}. Washington, DC: Catholic University of America Press, 2000. (Cited in Lesson 8 as a major modern scholarly exposition of Aquinas's metaphysical foundations.)
\end{itemize}
